\section*{Capítulo \ref{ch:b2:euler-bernoulli} --- Fluidos perfectos: Euler y Bernoulli}

\begin{solution}{exo:b2:euler-bernoulli:1}
$v_2 = 1.5 \times 4 = \qty{6.0}{m/s}$, $P_2 = 3.0 \times 10^5 - 500(36 - 2.25) =
\qty{2.83}{bar}$. A \qty{2}{cm}: $v = \qty{24}{m/s}$, $P = 3.0 \times 10^5 - 500
(576 - 2.25) = \qty{0.13}{bar}$ --- cerca de la presión de vapor. Un
resultado negativo significa que el caudal supuesto es imposible: el agua
cavita y el flujo queda estrangulado.
\end{solution}

\begin{solution}{exo:b2:euler-bernoulli:2}
$v = \sqrt{2 \times 15000/0.74} = \qty{201}{m/s}$; indicada, $\sqrt{2 \times 15000/
1.22} = \qty{157}{m/s}$. La velocidad indicada mide la \omterm{thm:b2:euler-bernoulli:bernoulli}{presión dinámica},
que es lo que siente el ala: la celeridad de entrada en pérdida es una
velocidad \emph{indicada} fija a cualquier altitud.
\end{solution}

\begin{solution}{exo:b2:euler-bernoulli:3}
Profundidad \qty{2.0}{m}: $v = \sqrt{2g \times 2} = \qty{6.3}{m/s}$; $D_V = 0.6 \times 10^{-4}
\times 6.3 = \qty{0.38}{L/s}$; tiempo de caída $\sqrt{2/g} = \qty{0.45}{s}$, alcance
\qty{2.8}{m}. El alcance $= 2\sqrt{h(H - h)}$ es simétrico en $h \leftrightarrow H -
h$: un orificio a \qty{1.0}{m} de profundidad (\qty{2}{m} sobre el suelo) llega al mismo
punto.
\end{solution}

\begin{solution}{exo:b2:euler-bernoulli:4}
$\Delta P = \tfrac12 \times 1.2 \times 900 = \qty{540}{Pa}$, con la presión exterior menor:
fuerza neta de \qty{54}{kN} \emph{hacia arriba}; el tejado pesa \qty{49}{kN}: se
levanta. (Los tejados reales fallan antes por los bordes, donde la corriente se desprende.)
\end{solution}

\begin{solution}{exo:b2:euler-bernoulli:5}
(a) $\tan\theta = 3/9.81$, $\theta = \qty{17}{\degree}$; $2 \times 0.306 = \qty{0.61}{m}$
(más alto atrás). (b) $\Delta P = \rho aL = \qty{6.0}{kPa}$. (c) $\tan\theta =
0.61$, diferencia \qty{1.22}{m}: la parte delantera sube \qty{0.61}{m} sobre la media de
\qty{0.40}{m}: \qty{1.01}{m} $>$ \qty{0.8}{m}, llega a la tapa. (d) $a = v^2/R =
\qty{5.6}{m/s^2}$: $\theta = \qty{30}{\degree}$, con el lado exterior más alto.
\end{solution}

\begin{solution}{exo:b2:euler-bernoulli:6}
(a) Paraboloide $z = z_0 + \Omega^2r^2/2g$, $\Omega = 4\pi = \qty{12.6}{rad/s}$. (b)
$\Omega^2R^2/2g = 158 \times 0.0225/19.6 = \qty{0.18}{m}$. (c) La altura media de un
paraboloide sobre el disco está a media altura: centro $0.20 - 0.09 = \qty{0.11}{m}$, borde
\qty{0.29}{m}. (d) El centro se seca cuando $\Omega^2R^2/4g = h_0$: $\Omega = \sqrt{4gh_0}/R =
\qty{18.7}{rad/s} = 3.0$ vueltas por segundo. (e) Un paraboloide enfoca los
rayos paralelos en un solo punto, a $f = g/2\Omega^2$; los telescopios de
espejo líquido hacen girar mercurio.
\end{solution}

\begin{solution}{exo:b2:euler-bernoulli:7}
(a) $v = \sqrt{2g \times 3} = \qty{7.7}{m/s}$; $D_V = \pi \times 10^{-4} \times 7.7 =
\qty{2.4}{L/s}$. (b) $P = P_0 - \rho g(2) - \tfrac12\rho v^2 = 100 - 19.6 - 29.4 =
\qty{51}{kPa}$. (c) $P_0 - \rho gh_c - \tfrac12\rho v^2 \ge \qty{2.3}{kPa}$: $h_c \le
\qty{7.0}{m}$. (d) El aire es compresible y ligero: no hay columna líquida
continua que transmita la presión, y el agua de ambos lados se queda quieta.
\end{solution}

\begin{solution}{exo:b2:euler-bernoulli:8}
(a) $\Delta P = mg/S = \qty{654}{Pa}$; $\tfrac12\rho v^2 = \qty{2.2}{kPa}$, $C_L = 0.30$.
(b) $v_{\text{extr}}^2 = 58^2 + 2 \times 654/1.22$: \qty{66.6}{m/s}. (c) $F_L' =
\qty{981}{N/m} = \rho v\Gamma$: $\Gamma = \qty{13.4}{m^2/s}$. (d) $v = \sqrt{2mg/
\rho SC_L} = \sqrt{19620/25.6} = \qty{28}{m/s}$.
\end{solution}

\begin{solution}{exo:b2:euler-bernoulli:9}
(a) $v_2 = \qty{2.5}{m/s}$; $\Delta P = \tfrac12 \times 1060 \times (6.25 - 0.25) =
\qty{3.2}{kPa}$, $P_2 = \qty{9.8}{kPa}$. (b) Sigue por encima de \qty{1}{kPa}: no. (c) A
\qty{0.05}{cm^2}, $v = \qty{5}{m/s}$ y $\Delta P = \qty{13}{kPa}$: $P_2 \approx 0$, por debajo
de la presión del tejido --- la pared colapsa, el flujo se detiene, la
presión se recupera y vuelve a abrirse: un aleteo (el soplo que oye el médico).
\end{solution}

\begin{solution}{exo:b2:euler-bernoulli:10}
(a) $c_PT_0 = c_PT + v^2/2$ con $c_P = \gamma R/(\gamma - 1)M_{\text{mol}}$ y
$c^2 = \gamma RT/M_{\text{mol}}$: $v^2/2c_PT = \tfrac12(\gamma - 1)M^2$. (b) $M =
0.85$: $T_0 = 220 \times 1.145 = \qty{252}{K}$; $M = 2$: $220 \times 1.8 = \qty{396}{K}$;
$M = 25$: $220 \times 126 = \qty{28000}{K}$ --- sin sentido: el aire se disocia
e ioniza y $c_P$ no es constante (las temperaturas de estancamiento reales
son de varios miles de kelvin). (c) $P_0/P = (1 + \tfrac{\gamma - 1}2M^2)^{\gamma/(\gamma
- 1)} \approx 1 + \tfrac{\gamma}2M^2 + \tfrac{\gamma}8M^4$, y $\gamma PM^2 = \rho v^2$:
$P_0 - P = \tfrac12\rho v^2(1 + M^2/4)$. Leer $v$ a partir de $\sqrt{2(P_0 - P)/\rho}$
lo sobreestima en $\sqrt{1 + M^2/4} - 1 \approx 9\%$ a $M = 0.85$.
\end{solution}

\begin{solution}{exo:b2:euler-bernoulli:11}
(a) Desde la superficie ($v \approx 0$, $P_0$, altura $h$) hasta la salida ($v$, $P_0$,
altura $0$), con $\int\partial_tv\,\dd s = \ell\,\dd v/\dd t$. (b) $\dd v/(V^2 -
v^2) = \dd t/2\ell$, $V^2 = 2gh$: $\operatorname{artanh}(v/V) = Vt/2\ell$. (c) $V =
\qty{9.9}{m/s}$; $\tanh x = 0.9$ en $x = 1.47$: $t = 2\ell x/V = \qty{15}{s}$. (d)
$\dd v/\dd t \approx \qty{-20}{m/s^2}$: $\Delta P = \rho\ell \times 20 = \qty{9.9}{bar}$.
El agua es ligeramente compresible: un cierre rápido lanza una onda de
presión, $\Delta P = \rho c\Delta v \approx \qty{140}{bar}$.
\end{solution}

\begin{solution}{exo:b2:euler-bernoulli:12}
(a) Estacionario, perfecto, incompresible e irrotacional: una única
constante para todo el flujo. (b) En la superficie $P = P_0$: $\tfrac12\rho v^2 + \rho gz =
0$, $z = -\Gamma^2/8\pi^2gr^2$. (c) Euler radial: $\partial_rP = \rho\Omega^2r$,
e hidrostática vertical: la superficie cumple $z(r) - z(a) = \Omega^2(r^2 -
a^2)/2g$. (d) Bañera: $\Omega = \Gamma/2\pi a^2 = \qty{191}{rad/s}$; $z(a) = -\Omega^2a^2/
2g = \qty{-4.7}{cm}$, y el núcleo añade otros \qty{4.7}{cm}: \qty{9.3}{cm} de profundidad.
Tornado: $\Delta P = \tfrac12\rho v_{\max}^2$ fuera más otro tanto dentro:
$\rho v_{\max}^2 = 1.2 \times 6400 = \qty{7.7}{kPa}$ por debajo de la ambiente en el centro.
\end{solution}

\begin{solution}{pb:b2:euler-bernoulli:1}
\textbf{1.} $\rho g \times 150 = \qty{14.7}{bar}$ sobre la atmosférica; $\tfrac12\rho gH_d^2
\times \qty{1}{m} = \qty{1.1e8}{N}$.

\textbf{2.} $\rho gH = \qty{19.6}{bar}$ manométricos, \qty{20.6}{bar} absolutos.

\textbf{3.} $v = 60/7.07 = \qty{8.5}{m/s}$; $\tfrac12\rho v^2 = \qty{0.36}{bar}$; $P =
P_0 + \rho gH - \tfrac12\rho v^2 = \qty{20.3}{bar}$ absolutos.

\textbf{4.} $v_j = \sqrt{2gH} = \qty{62.6}{m/s}$; $s = 60/62.6 = \qty{0.96}{m^2}$, $d =
\qty{1.1}{m}$.

\textbf{5.} $\tfrac12\rho D_Vv_j^2 = \tfrac12\rho D_V(2gH) = \rho gD_VH = \qty{118}{MW}$.

\textbf{6.} $m = \rho S\ell = \qty{2.8e6}{kg}$, $E_k = \tfrac12mv^2 = \qty{1.0e8}{J}$;
tránsito $\ell/v = \qty{47}{s}$.

\textbf{7.} $\rho gH = (P - P_0) + \tfrac12\rho v^2$: \qty{19.3}{bar} de presión
y \qty{0.36}{bar} de energía cinética; en la tobera la presión cae hasta
$P_0$ mientras $\tfrac12\rho v^2$ sube hasta $\rho gH$: justo después, $P = P_0$.

\textbf{8.} $\Delta h = \lambda(\ell/D)v^2/2g = 0.015 \times 133 \times 3.67 = \qty{7.3}{m}$,
el $3.7\%$ de la potencia.

\textbf{9.} $v = \qty{19.1}{m/s}$, $v^2/2g = \qty{18.6}{m}$, $\Delta h = 0.015 \times 200 \times
18.6 = \qty{56}{m}$: se pierde el $28\%$ --- el ahorro en acero se paga cada
segundo en energía.

\textbf{10.} $v_t = 60/\pi = \qty{19.1}{m/s}$; $\Delta P = 500(19.1^2 - 8.49^2) =
\qty{146}{kPa}$.

\textbf{11.} $\Delta P = (\rho_{\text{Hg}} - \rho)g\,\Delta h$: $\Delta h = 146000/(12600
\times 9.81) = \qty{1.18}{m}$.

\textbf{12.} $\Delta P \propto v^2 \propto D_V^2$: $D_V \propto \sqrt{\Delta h}$; la mitad
del caudal da la cuarta parte, \qty{0.30}{m}.

\textbf{13.} Una toma que sobresale se enfrenta a la corriente o la
perturba, y mide parte de la \omterm{thm:b2:euler-bernoulli:bernoulli}{presión dinámica} (es un Pitot), no la presión
estática.

\textbf{14.} $P = P_0 - \rho g \times 6 - \tfrac12\rho v^2 = 100 - 59 - 36 = \qty{5}{kPa}$:
cerca de la presión de vapor; el aire disuelto se libera y el agua puede
hervir (cavitación), rompiendo la columna. Hay que bajar la cota, o
instalar una ventosa y una bomba de vacío, o reducir el caudal.

\textbf{15.} $D_V/A = 60/2 \times 10^6 = \qty{3e-5}{m/s} = \qty{11}{cm/h}$.

\textbf{16.} $A\,\dd h/\dd t = -s\sqrt{2gh}$: $\sqrt h = \sqrt{h_0} - (s/A)\sqrt{g/2}\,t$.

\textbf{17.} $\Delta t = (A/s)\sqrt{2/g}(\sqrt{200} - \sqrt{190}) = 2.09 \times 10^6 \times
0.45 \times 0.36 = \qty{3.4e5}{s}$, es decir $3.9$ días; al caudal de diseño constante,
$10A/D_V = \qty{3.3e5}{s}$: casi lo mismo, pues la tobera se dimensionó
para el caudal de Torricelli con $h \approx H$.

\textbf{18.} $\ell\,\dd v/\dd t = gH - v^2/2$: $v = V\tanh(Vt/2\ell)$, $V = \qty{62.6}{m/s}$,
constante de tiempo $2\ell/V = \qty{13}{s}$.

\textbf{19.} En $v = V/2$: $\dd v/\dd t = (gH - V^2/8)/\ell = \qty{3.7}{m/s^2}$;
$P = P_0 + \rho gH - \tfrac12\rho v^2 - \rho\ell\,\dd v/\dd t = 1.0 + 19.6 - 4.9 - 14.7
= \qty{1.0}{bar}$: casi toda la carga está ocupada en acelerar la columna,
de modo que la presión en el fondo queda muy por debajo de su valor estacionario.

\textbf{20.} $|\dd v/\dd t| = 8.49/10 = \qty{0.85}{m/s^2}$: $\Delta P = \rho\ell \times 0.85 =
\qty{3.4}{bar}$.

\textbf{21.} $\Delta P = 1000 \times 1400 \times 8.49 = \qty{119}{bar}$, seis veces la
presión estática. El viaje de ida y vuelta de la onda dura $2\ell/c = \qty{0.57}{s}$:
un cierre más rápido que eso es ``súbito'' y recibe toda la sobrepresión
de Joukowsky; los cierres más lentos reciben menos.

\textbf{22.} $Av = A_s\,\dd z/\dd t$.

\textbf{23.} Euler a lo largo del túnel, desde el embalse (presión
$P_0 + \rho g\times$profundidad) hasta el pie del pozo (presión $P_0 + \rho g
(\text{profundidad} + z)$): $\rho L\,\dd v/\dd t = -\rho gz$; con $v = (A_s/A)\dot z$:
$\ddot z + (gA/LA_s)z = 0$; $T = 2\pi\sqrt{LA_s/gA} = 2\pi\sqrt{2000 \times 100/
68.7} = \qty{340}{s} \approx \qty{5.7}{min}$.

\textbf{24.} $\tfrac12\rho ALv^2 = \tfrac12\rho gA_sz_{\max}^2$ con $v = 60/7 =
\qty{8.6}{m/s}$: $z_{\max} = v\sqrt{AL/gA_s} = \qty{32}{m}$.

\textbf{25.} Presa: hidrostática, $\rho gh$. Turbina con las válvulas
cerradas: lo mismo, $\rho gH$. Tubería forzada en servicio: Bernoulli,
$\tfrac12\rho v^2$ tomado de la carga estática. Golpe de ariete:
compresibilidad, $\rho c\Delta v$. Chimenea: Euler no estacionario a lo
largo del túnel, una oscilación de periodo $2\pi\sqrt{LA_s/gA}$.
\end{solution}
